% Options for packages loaded elsewhere
\PassOptionsToPackage{unicode}{hyperref}
\PassOptionsToPackage{hyphens}{url}
\PassOptionsToPackage{dvipsnames,svgnames,x11names}{xcolor}
%
\documentclass[
  11pt,
]{article}
\usepackage{amsmath,amssymb}
\usepackage{iftex}
\ifPDFTeX
  \usepackage[T1]{fontenc}
  \usepackage[utf8]{inputenc}
  \usepackage{textcomp} % provide euro and other symbols
\else % if luatex or xetex
  \usepackage{unicode-math} % this also loads fontspec
  \defaultfontfeatures{Scale=MatchLowercase}
  \defaultfontfeatures[\rmfamily]{Ligatures=TeX,Scale=1}
\fi
\usepackage{lmodern}
\ifPDFTeX\else
  % xetex/luatex font selection
\fi
% Use upquote if available, for straight quotes in verbatim environments
\IfFileExists{upquote.sty}{\usepackage{upquote}}{}
\IfFileExists{microtype.sty}{% use microtype if available
  \usepackage[]{microtype}
  \UseMicrotypeSet[protrusion]{basicmath} % disable protrusion for tt fonts
}{}
\makeatletter
\@ifundefined{KOMAClassName}{% if non-KOMA class
  \IfFileExists{parskip.sty}{%
    \usepackage{parskip}
  }{% else
    \setlength{\parindent}{0pt}
    \setlength{\parskip}{6pt plus 2pt minus 1pt}}
}{% if KOMA class
  \KOMAoptions{parskip=half}}
\makeatother
\usepackage{xcolor}
\usepackage[margin=1in]{geometry}
\usepackage{longtable,booktabs,array}
\usepackage{calc} % for calculating minipage widths
% Correct order of tables after \paragraph or \subparagraph
\usepackage{etoolbox}
\makeatletter
\patchcmd\longtable{\par}{\if@noskipsec\mbox{}\fi\par}{}{}
\makeatother
% Allow footnotes in longtable head/foot
\IfFileExists{footnotehyper.sty}{\usepackage{footnotehyper}}{\usepackage{footnote}}
\makesavenoteenv{longtable}
\usepackage{graphicx}
\makeatletter
\def\maxwidth{\ifdim\Gin@nat@width>\linewidth\linewidth\else\Gin@nat@width\fi}
\def\maxheight{\ifdim\Gin@nat@height>\textheight\textheight\else\Gin@nat@height\fi}
\makeatother
% Scale images if necessary, so that they will not overflow the page
% margins by default, and it is still possible to overwrite the defaults
% using explicit options in \includegraphics[width, height, ...]{}
\setkeys{Gin}{width=\maxwidth,height=\maxheight,keepaspectratio}
% Set default figure placement to htbp
\makeatletter
\def\fps@figure{htbp}
\makeatother
\setlength{\emergencystretch}{3em} % prevent overfull lines
\providecommand{\tightlist}{%
  \setlength{\itemsep}{0pt}\setlength{\parskip}{0pt}}
\setcounter{secnumdepth}{-\maxdimen} % remove section numbering
\ifLuaTeX
  \usepackage{selnolig}  % disable illegal ligatures
\fi
\IfFileExists{bookmark.sty}{\usepackage{bookmark}}{\usepackage{hyperref}}
\IfFileExists{xurl.sty}{\usepackage{xurl}}{} % add URL line breaks if available
\urlstyle{same}
\hypersetup{
  pdftitle={Design of a Reproducible Healthy Aging Policy Observatory: Infrastructure for Trustworthy Aging-Policy Evidence in Canada},
  pdfauthor={Quangui Huang (Giorgio), Healthy Aging Intelligence Lab (HAIL)},
  colorlinks=true,
  linkcolor={RoyalBlue},
  filecolor={Maroon},
  citecolor={Blue},
  urlcolor={RoyalBlue},
  pdfcreator={LaTeX via pandoc}}

\title{Design of a Reproducible Healthy Aging Policy Observatory:
Infrastructure for Trustworthy Aging-Policy Evidence in Canada}
\author{Quangui Huang (Giorgio), Healthy Aging Intelligence Lab (HAIL)}
\date{Preprint --- draft, not peer reviewed}

\begin{document}
\maketitle

\hypertarget{abstract}{%
\subsection{Abstract}\label{abstract}}

The evidence base for aging policy suffers less from a shortage of
analysis than from a shortage of \emph{trustworthy infrastructure}:
policy text and outcome data live apart, published numbers rarely trace
to their sources, and descriptive co-movement is routinely dressed as
causation. \textbf{We argue these are infrastructure problems, not
analysis problems} --- and that the response is a reproducible research
instrument, an \emph{observatory}, rather than another dashboard. We
frame the design around a conceptual model, the \textbf{Healthy Aging
Policy Intelligence Cycle} (Observation → Evidence → Indicator → Policy
→ Outcome → Feedback) --- a policy-level instance of a Learning Health
System --- and ask four design-evaluation questions: can every published
number be traced to its exact source (provenance); does the pipeline
reproduce identical results (reproducibility); can the AI layer emit an
ungrounded citation (grounding); and are composite comparisons robust to
method choices (robustness). We present the Canadian Healthy Aging
Policy Observatory, whose data model makes provenance structural
(immutable observations bound to versioned, content-addressed dataset
snapshots) and whose analytics enforce an association-versus-causation
distinction in software. Instantiated on Nova Scotia and the federal
level with a Care Access focus, the instrument answers the design
questions affirmatively: 100\% of published values resolve to an
upstream source, the pipeline reproduces deterministically from
committed inputs, the assistant cannot cite outside a fixed evidence
pack, and the headline index ordering is invariant across three
weighting schemes (composite spread 2.9 points). \textbf{The design
itself is the contribution}: a reproducible, provenance-first
observatory is a durable paradigm for aging-policy evidence that
compounds across a multi-paper program, in which the index methodology,
the grounded-AI framework, and the causal-evaluation studies are each
deepened in their own right.

\textbf{Keywords:} healthy aging, health informatics, learning health
system, policy evaluation, reproducible research infrastructure, data
provenance, composite indicators, interrupted time series,
retrieval-augmented generation, Nova Scotia.

\begin{center}\rule{0.5\linewidth}{0.5pt}\end{center}

\hypertarget{1-introduction}{%
\subsection{1. Introduction}\label{1-introduction}}

\hypertarget{11-the-problem-is-infrastructure-not-analysis}{%
\subsubsection{1.1 The problem is infrastructure, not
analysis}\label{11-the-problem-is-infrastructure-not-analysis}}

Canada, like most high-income countries, is aging quickly, and
governments at every level respond with policy --- home-care investment,
long-term-care (LTC) capacity, dementia strategies,
seniors\textquotesingle{} income supports, digital-inclusion programs.
Yet the machinery for \emph{evaluating} those policies lags the
machinery for \emph{announcing} them. Three gaps recur:

\begin{enumerate}
\def\labelenumi{\arabic{enumi}.}
\tightlist
\item
  \textbf{Fragmentation.} Policy text, budgets, and target populations
  live in press releases and PDFs; outcome data live in separate
  agencies (Statistics Canada, the Canadian Institute for Health
  Information, provincial open-data portals). Nothing links a policy to
  the indicators it was meant to move.
\item
  \textbf{Irreproducibility.} Published dashboards and reports rarely
  let a reader trace a headline number to the specific upstream table,
  vintage, and transformation that produced it; re-running an analysis a
  year later often yields a different number for unstated reasons.
\item
  \textbf{Causal overreach.} Descriptive co-movement ("spending rose and
  ED visits fell") is routinely presented in language implying
  causation, without the design or assumptions that would justify it.
\end{enumerate}

Our central claim is that \textbf{these are infrastructure problems, not
analysis problems.} No new estimator fixes a number that cannot be
traced to its source, an analysis that cannot be reproduced, or a claim
whose causal status is ambiguous by default. What is missing is a
standing, reproducible \emph{instrument} --- an \textbf{observatory} ---
that continuously ingests policies and outcomes, quantifies them with a
documented and versioned method, and makes honest inference the path of
least resistance. The distinction is the one astronomy draws between a
telescope and a press photo: the value is in the
instrument\textquotesingle s reliability and in the fact that others can
point it themselves.

\hypertarget{12-why-an-aging-policy-observatory-is-health-informatics}{%
\subsubsection{1.2 Why an aging-policy observatory is health
informatics}\label{12-why-an-aging-policy-observatory-is-health-informatics}}

Policy is an upstream determinant of health: home-care funding, LTC
staffing, and income supports shape the functional ability,
independence, and outcomes of older adults as surely as any bedside
intervention. Health informatics has correspondingly broadened from
clinical decision support toward \textbf{population- and policy-level}
decision support and \textbf{knowledge translation} --- turning evidence
into actionable, traceable syntheses for decision-makers. The
observatory is precisely such an instrument, and its design mirrors a
core health-informatics idea: the \textbf{Learning Health System} (LHS)
{[}11{]}, in which data generate knowledge that improves practice and
thereby generates new data, in a continuous loop. Where a clinical LHS
closes that loop around care delivery, the observatory closes it around
\emph{policy} --- an LHS for healthy-aging governance. Reading the
design through evidence, knowledge translation, and decision support is
what makes it a health-informatics contribution rather than a
software-engineering one.

\hypertarget{13-research-questions}{%
\subsubsection{1.3 Research questions}\label{13-research-questions}}

Because the contribution is a design, we evaluate it as a design. We ask
four \textbf{design-evaluation questions} whose answers are properties
of the artifact, not outcomes of a clinical trial:

\begin{itemize}
\tightlist
\item
  \textbf{RQ1 (Provenance).} Can every published quantitative value be
  traced to the exact upstream source it came from?
\item
  \textbf{RQ2 (Reproducibility).} Does the pipeline reproduce identical
  results from the same inputs, and is ingestion idempotent under
  unchanged upstream data?
\item
  \textbf{RQ3 (Grounding).} Can the AI layer emit a factual claim or
  citation that is not grounded in a fixed, visible evidence set?
\item
  \textbf{RQ4 (Robustness).} Are the composite index\textquotesingle s
  headline comparisons robust to its most contestable method choice, the
  domain weighting?
\end{itemize}

§7 answers each with evidence from the instantiated instrument.

\hypertarget{14-contribution}{%
\subsubsection{1.4 Contribution}\label{14-contribution}}

The primary contribution of this paper is \textbf{the design of a
reproducible healthy-aging policy observatory, together with the
conceptual framework it realizes} --- the Healthy Aging Policy
Intelligence Cycle (§3). The design\textquotesingle s defining
commitments are provenance as a structural property, reproducibility by
construction, and an association/causation discipline enforced in
software. The observatory \emph{instantiates} three modules --- a
documented composite index (HAPI), a quasi-experimental analytics layer,
and a grounded AI research assistant --- but these are presented here as
evidence that the framework runs end-to-end, not as the
paper\textquotesingle s contribution; each is the subject of its own
paper in the program (§8). This scoping is deliberate: a single,
well-evaluated design claim is worth more than an inventory of
subsystems.

\hypertarget{2-related-work}{%
\subsection{2. Related work}\label{2-related-work}}

\textbf{Learning health systems and health-informatics infrastructure.}
The LHS concept {[}11{]} frames continuous improvement as a
data→knowledge→practice→data cycle; most LHS work targets clinical care.
We extend the idea \emph{upstream} to policy, and we take seriously the
informatics infrastructure --- provenance, versioning, reproducibility
--- that an LHS presupposes but often under-specifies. Our
data-stewardship commitments follow the FAIR principles {[}9{]}.

\textbf{Composite indices of aging.} The WHO \emph{World Report on
Ageing and Health} {[}1{]} and the \emph{Decade of Healthy Ageing}
{[}2{]} articulate intrinsic capacity and functional ability; the
European Active Ageing Index {[}5{]} and HelpAge\textquotesingle s
Global AgeWatch Index {[}4{]} operationalize domains into composite
scores. HAPI is in this tradition but is policy-facing, reproducible by
construction, and transparent about its own weighting; its methodology
is developed in full in Paper 2 (§8) and appears here only as an
instantiated module. We follow OECD/JRC guidance {[}3{]} on transparency
and sensitivity analysis.

\textbf{Policy monitors and quasi-experimental evaluation.}
Government/NGO trackers catalog policies but seldom integrate outcomes
or support inference; academic evaluations produce rigorous
single-policy studies but rarely as reusable infrastructure. The
segmented-regression interrupted-time-series literature {[}6, 7, 8{]}
provides the reference designs; the observatory operationalizes them
under an explicit tagging discipline and defers their empirical
application to Paper 4.

\textbf{LLMs for policy analysis.} Retrieval-augmented generation (RAG)
{[}10{]} makes it feasible to draft grounded syntheses over a fixed
corpus; the risk is fabricated citations and unsupported claims. Our
assistant constrains the model to cite \emph{only} items in a
pre-assembled evidence pack, making fabrication structurally detectable;
the framework is developed in Paper 3.

\hypertarget{3-conceptual-framework-the-healthy-aging-policy-intelligence-cycle}{%
\subsection{3. Conceptual framework: the Healthy Aging Policy
Intelligence
Cycle}\label{3-conceptual-framework-the-healthy-aging-policy-intelligence-cycle}}

The design is organized by a conceptual model that doubles as its theory
of value. Healthy-aging policy intelligence is a cycle of six stages
(Figure 1):

\includegraphics{figures/intelligence-cycle.png}

\begin{quote}
\textbf{Figure 1.} The Healthy Aging Policy Intelligence Cycle. Six
stages form a closed loop; the italic labels mark the four points where
the design hardens trust (traceability, reproducibility, causal honesty,
and grounding).
\end{quote}

\begin{itemize}
\tightlist
\item
  \textbf{Observation} --- an immutable, source-linked datum ingested
  from an upstream table.
\item
  \textbf{Evidence} --- observations assembled and quality-checked into
  an interpretable body, with literature and prior findings.
\item
  \textbf{Indicator} --- evidence normalized into a documented,
  comparable measure (HAPI).
\item
  \textbf{Policy} --- a versioned policy record linked to the indicators
  it intends to move.
\item
  \textbf{Outcome} --- the movement of indicators over time, analyzed
  under an association/causation discipline.
\item
  \textbf{Feedback} --- findings and syntheses returned to researchers
  and policymakers, which motivate new observations, closing the loop.
\end{itemize}

Two properties make this more than a diagram. First, it is a
\textbf{policy-level Learning Health System} {[}11{]}: each turn of the
cycle converts data into knowledge that informs policy and thereby
generates new data. Second, it names exactly the points where trust is
usually lost --- the arrows --- and the design hardens each:
Observation→Evidence is made \emph{traceable} (lineage), Indicator is
made \emph{reproducible and version-stamped}, Outcome is made
\emph{honest} (causal tagging), and Feedback is made \emph{grounded}
(cited AI). The five platform modules (§4) are the machinery that turns
this cycle; the rest of the paper shows the machinery exists, runs, and
can be trusted.

\hypertarget{4-system-design}{%
\subsection{4. System design}\label{4-system-design}}

\hypertarget{41-modules-as-cycle-stages}{%
\subsubsection{4.1 Modules as cycle
stages}\label{41-modules-as-cycle-stages}}

The Observatory is five modules over one database and one provenance
discipline (Figure 2), each realizing part of the cycle:

\begin{enumerate}
\def\labelenumi{\arabic{enumi}.}
\tightlist
\item
  \textbf{Policy Library} (Policy) --- a jurisdiction-aware,
  time-stamped, versioned record of aging policies: department, budget,
  target population, lifecycle, themes, and AI-assisted summaries.
\item
  \textbf{Data Hub} (Observation → Evidence) --- connectors that ingest
  open data into immutable observations with full lineage and automated
  quality checks.
\item
  \textbf{Indicators / HAPI} (Indicator) --- a documented composite over
  six domains, computed from observations and stored with a method
  version.
\item
  \textbf{Policy Analytics} (Outcome) --- descriptive trends and
  quasi-experimental designs under an explicit association/causation
  tag.
\item
  \textbf{AI Research Assistant} (Feedback) --- retrieval over the
  platform\textquotesingle s own stores that produces cited evidence
  packs and draft reviews.
\end{enumerate}

\includegraphics{figures/system-context.png}

\begin{quote}
\textbf{Figure 2.} System context: external public sources flow into the
five Observatory modules, which produce papers, dashboards, and
literature reviews. (Source:
\href{../../docs/01-platform-overview.md}{\texttt{docs/01-platform-overview.md}}.)
\end{quote}

\hypertarget{42-architecture}{%
\subsubsection{4.2 Architecture}\label{42-architecture}}

The platform is a monorepo (Figure 3): a \textbf{Next.js/TypeScript} web
app (\texttt{apps/web}) for the researcher-facing surfaces; a
\textbf{Python} pipeline (\texttt{pipeline/hapi\_pipeline}) for
ingestion, scoring, and analysis; a \textbf{PostgreSQL} schema with
forward-only migrations (\texttt{db/}); and a shared contracts package.
The pipeline is the reproducible core (deterministic transforms,
versioned method); the web app is a read-mostly view over what the
pipeline produced. A number in the UI is never computed in the browser
--- it renders a stored, lineage-bearing record. The AI layer wraps a
hosted LLM behind narrow, audited entry points, and every generation is
logged.

\includegraphics{figures/architecture.png}

\begin{quote}
\textbf{Figure 3.} Architecture: the Python pipeline writes to
PostgreSQL (raw → transformed → indicators → analytics); the Next.js app
reads those tables and invokes the AI layer. (Source:
\href{../../docs/02-architecture.md}{\texttt{docs/02-architecture.md}}.)
\end{quote}

\hypertarget{43-design-principles}{%
\subsubsection{4.3 Design principles}\label{43-design-principles}}

\emph{Foundation before breadth} (get the model, lineage, and method
right on a narrow slice, then widen additively); \emph{every artifact is
demonstrable and reproducible}; and \emph{researcher-first}
(public/decision-support surfaces follow a credible research core).

\hypertarget{5-data-model-making-provenance-structural}{%
\subsection{5. Data model: making provenance
structural}\label{5-data-model-making-provenance-structural}}

The schema (Figure 4) centers on entities that make provenance and
reproducibility \emph{structural} rather than conventional:

\begin{itemize}
\tightlist
\item
  \textbf{Jurisdiction} --- a hierarchical tree (Canada → Federal / Nova
  Scotia → \ldots); adding a province is a row, not a schema change.
\item
  \textbf{Policy / policy\_version} --- versioned policy records linked
  to a jurisdiction.
\item
  \textbf{Indicator} --- the definition of a measure: meaning, unit,
  normalization, per-capita rule.
\item
  \textbf{DataSource} --- upstream provider metadata (agency, licence,
  cadence, access method).
\item
  \textbf{DatasetVersion} --- a specific \emph{retrieval}: timestamp +
  content \textbf{checksum}, the unit of idempotency.
\item
  \textbf{Observation} --- an immutable (jurisdiction × indicator ×
  time) value bound to its DatasetVersion; never mutated --- a new
  vintage yields a new version and new observations.
\item
  \textbf{hapi\_score} --- a computed domain/composite score at a given
  \texttt{method\_version}.
\item
  \textbf{policy\_indicator, indicator\_source} --- the many-to-many
  links that make the platform a small knowledge graph.
\item
  \textbf{analysis\_finding, literature} --- analytic results
  (tier-tagged) and the literature knowledge base.
\end{itemize}

The key property is the \textbf{lineage chain}: every value resolves as
\texttt{Observation\ →\ DatasetVersion\ →\ DataSource}. Because
observations are immutable and dataset versions are content-addressed,
analyses are reproducible and ingestion is idempotent (re-ingesting
unchanged upstream data is a no-op). This is what lets the platform
claim any published number is auditable to the exact upstream table it
came from --- and it is what §7 evaluates directly (RQ1, RQ2).

\includegraphics{figures/data-model-er.png}

\begin{quote}
\textbf{Figure 4.} Entity--relationship model. Every value resolves
along the lineage chain Observation → DatasetVersion → DataSource;
indicators roll up into HAPI scores. (Source:
\href{../../docs/03-data-model.md}{\texttt{docs/03-data-model.md}}.)
\end{quote}

\hypertarget{6-instantiating-the-cycle}{%
\subsection{6. Instantiating the
cycle}\label{6-instantiating-the-cycle}}

To show the framework runs end-to-end, the Observatory instantiates its
three analytical stages. We summarize each here as evidence; the full
methodology of each is the subject of its own paper (§8).

\textbf{Indicator --- HAPI.} The Healthy Aging Policy Index scores how
well a jurisdiction supports healthy aging (0--100) across six domains
(Health, Independence, Social Participation, Financial Security, Care
Access, Digital Inclusion). Each indicator declares a direction, a
min-max normalization against a documented reference range, and, for
counts, a per-capita rule; the composite is a weighted mean over
available domains, renormalized by the domains present. Crucially,
weighting --- the most contestable step --- is treated as an auditable
object: three schemes (equal; a theory-anchored \emph{expert} tier from
the WHO {[}1{]} and AgeWatch {[}4{]} frameworks; and an
\emph{empirical}, coefficient-of-variation scheme) are reported side by
side per OECD/JRC guidance {[}3{]}, and every score carries a
\texttt{method\_version}. \emph{The index\textquotesingle s methodology,
validation, and sensitivity analysis are developed fully in Paper 2.}

\textbf{Outcome --- analytics under a causal discipline.} Results are
two-tier and always labelled. \textbf{Tier 1 (Association)}: trends,
policy-event overlays, correlations. \textbf{Tier 2 (Causal)}:
quasi-experimental designs --- interrupted time series (ITS),
difference-in-differences, synthetic control --- each carrying its
assumptions as attached metadata. The reference Tier-2 design is
segmented-regression ITS {[}7, 8{]} with Newey--West (HAC) inference
{[}6{]}; it refuses to report coefficients below three points per
segment, returning an explicit \texttt{insufficient\_data} status
instead. A finding is tagged \texttt{Causal} \emph{only} together with
its named assumptions, never upgraded silently from an association.
\emph{The empirical evaluations are the subject of Paper 4.}

\textbf{Feedback --- grounded AI assistant.} Given a topic, the
assistant retrieves an \textbf{evidence pack} (policies \texttt{P\#},
literature \texttt{L\#}, findings \texttt{F\#}) from the
platform\textquotesingle s own stores, then drafts a review \emph{from
that exact pack} under a contract: every claim ends with a citation tag
resolving to a pack item; only the pack may be used; analytic tiers are
preserved, never upgraded; the output is a draft for a human, not a
verdict. Because the pack is fixed and visible, a fabricated citation is
structurally detectable, and every generation is logged. \emph{The
grounded-AI framework is developed in Paper 3.}

\hypertarget{7-design-evaluation}{%
\subsection{7. Design evaluation}\label{7-design-evaluation}}

We evaluate the design against the four research questions, using the
instrument instantiated on Nova Scotia (CA-NS) and the federal level
(CA) with a Care Access focus. Unless noted, results reproduce
deterministically from committed fixtures (captured from real upstream
tables) via \texttt{ingest\ →\ score\ →\ analyze\ →\ paper-tables}; the
production deployment runs the same connectors against live sources. The
instantiated corpus:

\begin{longtable}[]{@{}lr@{}}
\toprule\noalign{}
Asset & Count \\
\midrule\noalign{}
\endhead
\bottomrule\noalign{}
\endlastfoot
Policies & 26 \\
Indicators (distinct, observed) & 34 \\
Observations & 231 \\
Dataset versions & 17 \\
Data sources & 17 \\
HAPI scores & 100 \\
Analytic findings & 3 \\
Literature references & 7 \\
Jurisdictions & 3 \\
\end{longtable}

\textbf{RQ1 --- Provenance: can every number be traced?} Yes. Every one
of the 231 observations is bound to a DatasetVersion (17) and thence a
DataSource (17); every HAPI score and figure is computed from those
observations; and the Data Hub renders the full
\texttt{Observation\ →\ DatasetVersion\ →\ DataSource} chain for each
value. Traceability is not a sampled audit but a structural invariant of
the schema (§5): a value that lacks lineage cannot be stored.
\emph{Result: 100\% of published quantitative values are
source-traceable.}

\textbf{RQ2 --- Reproducibility: does the pipeline reproduce,
idempotently?} Yes. The transforms are deterministic, so the same inputs
yield identical scores and findings; re-ingesting unchanged upstream
data is a no-op because dataset versions are content-addressed by
checksum. The evaluation corpus above regenerates from committed
fixtures with a single command, which is how the numbers in this paper
were produced. \emph{Result: deterministic reproduction; idempotent
ingestion.}

\textbf{RQ3 --- Grounding: can the AI emit an ungrounded citation?} No,
by construction. The assistant may cite only ids present in the
pre-assembled evidence pack; any tag that does not resolve to a pack
item is detectable as fabrication, and every generation is logged for
audit. Grounding is therefore a property of the retrieval-then-constrain
design rather than of model behavior. \emph{Result: zero ungrounded
citations are representable; violations are structurally detectable.}

\textbf{RQ4 --- Robustness: are index comparisons robust to weighting?}
Yes. Under the three schemes, the composite for each jurisdiction at its
latest period is:

\begin{longtable}[]{@{}llrrr@{}}
\toprule\noalign{}
Jurisdiction & Period & equal & expert & empirical \\
\midrule\noalign{}
\endhead
\bottomrule\noalign{}
\endlastfoot
CA (federal) & 2024 & 43.2 & 41.7 & 44.6 \\
CA-NS (Nova Scotia) & 2026 & 50.1 & 50.1 & 50.1 \\
\end{longtable}

The \textbf{maximum composite spread across schemes is 2.9 points}, and
the NS-above-federal ordering holds under all three weightings --- so
the headline comparison does not hinge on the weighting choice. For
completeness, the computed v1 domain profiles (latest available score
per domain) are:

\begin{longtable}[]{@{}lrrrrrrr@{}}
\toprule\noalign{}
Jurisdiction & Overall & Health & Independence & Social Part. &
Financial Sec. & Care Access & Digital Incl. \\
\midrule\noalign{}
\endhead
\bottomrule\noalign{}
\endlastfoot
CA (federal) & 42 & 48 & 46 & 50 & 56 & 28 & 72 \\
CA-NS (Nova Scotia) & 50 & 38 & 44 & 69 & 21 & 50 & 59 \\
\end{longtable}

\textbf{Causal discipline in action.} The three ITS designs illustrate
the Outcome-stage guarantee: two attached to real NS policies return
\texttt{insufficient\_data} (1/3 and 5/2 points per segment --- below
the threshold), so no coefficients are asserted; a third,
adequately-covered illustrative series (5/5) yields a significant
downward slope change (−715.7; 95\% CI −1125.1..−306.3; p \textless{}
0.001; R² = 0.934), reported as \textbf{Causal(ITS)} only alongside its
four assumptions. That the platform \emph{demonstrates a design without
estimating an effect on thin data} is the honesty property, evaluated
live.

Together these answer the design questions affirmatively: the
observatory is traceable, reproducible, grounded, and robust --- the
properties the Intelligence Cycle (§3) identifies as where trust is won
or lost.

\hypertarget{8-the-research-program}{%
\subsection{8. The research program}\label{8-the-research-program}}

Scoping Paper 1 to the observatory design lets each analytical module be
developed as a focused study that both deepens and reuses the platform:

\begin{longtable}[]{@{}lll@{}}
\toprule\noalign{}
Paper & Contribution & Instantiates \\
\midrule\noalign{}
\endhead
\bottomrule\noalign{}
\endlastfoot
\textbf{1 (this paper)} & Design of a reproducible healthy-aging policy
observatory + the Policy Intelligence Cycle & Research infrastructure \\
\textbf{2} & The HAPI methodology --- a transparent, sensitivity-tested
composite-indicator framework & Indicator stage \\
\textbf{3} & AI-assisted policy intelligence --- grounded retrieval +
evidence synthesis & Feedback stage \\
\textbf{4} & Causal evaluation of Canadian aging policies --- ITS / DiD
/ synthetic control & Outcome stage \\
\end{longtable}

Each paper deposits assets back into the platform (a curated corpus, a
validated index, reusable agent tooling, worked evaluations), so the
program compounds rather than restarts --- the strategic point of
building an instrument rather than a study.

\hypertarget{9-limitations-and-ethics}{%
\subsection{9. Limitations and ethics}\label{9-limitations-and-ethics}}

\textbf{Limitations.} (1) \emph{Coverage.} With two jurisdictions, the
empirical weighting and any cross-jurisdiction inference are indicative;
the design\textquotesingle s value is in the method and its additive
extensibility. (2) \emph{Normalization ranges are choices.} Min-max
reference ranges are documented and contestable; RQ4\textquotesingle s
robustness check mitigates but does not eliminate this --- a composite
index is a lens, not a measurement. (3) \emph{Causal identification.}
ITS and related designs rest on assumptions (no coincident shocks,
correct functional form) that real timelines strain; thin data may
permit only \emph{demonstrating} a design, which the platform reports
honestly. (4) \emph{Design evaluation, not outcome evaluation.} We
evaluate whether the instrument is trustworthy (RQ1--RQ4), not whether
any policy improved health --- the latter is Paper 4\textquotesingle s
task.

\textbf{Ethics.} The platform uses aggregate, published statistics, not
individual records, and respects Canada\textquotesingle s federal
privacy law (PIPEDA) as a baseline. Most importantly, the design encodes
an ethic of \emph{epistemic honesty}: association is never dressed as
causation, index choices are exposed for scrutiny, and AI output is
grounded and cited by construction. For infrastructure that may inform
decisions about a vulnerable population, that honesty is a safety
property.

\hypertarget{10-conclusion}{%
\subsection{10. Conclusion}\label{10-conclusion}}

We have argued that the weaknesses of aging-policy evidence are
infrastructure problems, not analysis problems, and presented the design
of a reproducible observatory that treats them as such --- organized by
the Healthy Aging Policy Intelligence Cycle, a policy-level Learning
Health System. Evaluated as a design, the instrument is traceable (100\%
of values source-linked), reproducible (deterministic, idempotent),
grounded (no representable ungrounded citation), and robust (index
ordering invariant to weighting, spread 2.9 points). The contribution is
the design itself: a durable paradigm for trustworthy aging-policy
evidence, on which a focused program of methodology, grounded-AI, and
causal-evaluation papers builds.

\hypertarget{acknowledgements}{%
\subsection{Acknowledgements}\label{acknowledgements}}

This work uses publicly available data from Statistics Canada, the
Canadian Institute for Health Information (CIHI), and the Nova Scotia
Open Data portal; we thank these providers for open access to the
underlying tables. The author received no specific funding for this
work.

\hypertarget{data-and-code-availability}{%
\subsection{Data and code
availability}\label{data-and-code-availability}}

The platform --- schema, pipeline, methodology, and web application ---
is developed in the open at
\url{https://github.com/GiorgioHuang/aging-policy-lab}. Every
quantitative result reproduces from the pipeline commands in the
paper\textquotesingle s repository README.

\hypertarget{11-references}{%
\subsection{11. References}\label{11-references}}

\emph{Aging frameworks and composite indices.}

\begin{enumerate}
\def\labelenumi{\arabic{enumi}.}
\tightlist
\item
  World Health Organization. \emph{World Report on Ageing and Health.}
  Geneva: WHO, 2015. ISBN 978-92-4-156504-2.
\item
  World Health Organization. \emph{Decade of Healthy Ageing 2021--2030.}
  Geneva: WHO, 2020.
\item
  Nardo, M., Saisana, M., Saltelli, A., Tarantola, S., Hoffmann, A., \&
  Giovannini, E. \emph{Handbook on Constructing Composite Indicators:
  Methodology and User Guide.} Paris: OECD Publishing, 2008.
  doi:10.1787/9789264043466-en.
\item
  HelpAge International. \emph{Global AgeWatch Index 2015: Insight
  Report.} London: HelpAge International, 2015.
\item
  Zaidi, A., Gasior, K., Hofmarcher, M. M., Lelkes, O., Marin, B.,
  Rodrigues, R., Schmidt, A., Vanhuysse, P., \& Zolyomi, E. \emph{Active
  Ageing Index 2012: Concept, Methodology and Final Results.} Vienna:
  European Centre, 2013.
\end{enumerate}

\emph{Quasi-experimental methods (interrupted time series).}

\begin{enumerate}
\def\labelenumi{\arabic{enumi}.}
\setcounter{enumi}{5}
\tightlist
\item
  Newey, W. K., \& West, K. D. "A Simple, Positive Semi-Definite,
  Heteroskedasticity and Autocorrelation Consistent Covariance Matrix."
  \emph{Econometrica}, 55(3), 703--708, 1987. doi:10.2307/1913610.
\item
  Wagner, A. K., Soumerai, S. B., Zhang, F., \& Ross-Degnan, D.
  "Segmented regression analysis of interrupted time series studies in
  medication use research." \emph{Journal of Clinical Pharmacy and
  Therapeutics}, 27(4), 299--309, 2002.
  doi:10.1046/j.1365-2710.2002.00430.x.
\item
  Lopez Bernal, J., Cummins, S., \& Gasparrini, A. "Interrupted time
  series regression for the evaluation of public health interventions: a
  tutorial." \emph{International Journal of Epidemiology}, 46(1),
  348--355, 2017. doi:10.1093/ije/dyw098.
\end{enumerate}

\emph{Learning health systems, reproducible infrastructure, and grounded
generation.}

\begin{enumerate}
\def\labelenumi{\arabic{enumi}.}
\setcounter{enumi}{8}
\tightlist
\item
  Wilkinson, M. D., Dumontier, M., Aalbersberg, Ij. J., et al. "The FAIR
  Guiding Principles for scientific data management and stewardship."
  \emph{Scientific Data}, 3, 160018, 2016. doi:10.1038/sdata.2016.18.
\item
  Lewis, P., Perez, E., Piktus, A., Petroni, F., Karpukhin, V., Goyal,
  N., Küttler, H., Lewis, M., Yih, W., Rocktäschel, T., Riedel, S., \&
  Kiela, D. "Retrieval-Augmented Generation for Knowledge-Intensive NLP
  Tasks." \emph{NeurIPS} 33, 2020. arXiv:2005.11401.
\item
  Friedman, C. P., Wong, A. K., \& Blumenthal, D. "Achieving a
  Nationwide Learning Health System." \emph{Science Translational
  Medicine}, 2(57), 57cm29, 2010. doi:10.1126/scitranslmed.3001456.
\end{enumerate}

\emph{Primary data sources.}

\begin{enumerate}
\def\labelenumi{\arabic{enumi}.}
\setcounter{enumi}{11}
\tightlist
\item
  Statistics Canada. Data accessed via the Web Data Service (WDS).
  Tables used include 13-10-0971-01 (health-adjusted life expectancy),
  13-10-0096-01 (Canadian Community Health Survey --- health
  characteristics), 13-10-0374-01 (Canadian Survey on Disability),
  13-10-0789-01 (Canadian Health Survey on Seniors --- ADL),
  11-10-0135-01 (low-income statistics), and 14-10-0202-01 (employment
  by industry, SEPH). \emph{{[}TODO: confirm each
  table\textquotesingle s exact vintage as cited in the Data Hub.{]}}
\item
  Canadian Institute for Health Information (CIHI). Integrated interRAI
  Reporting System (IRRS) and long-term-care data holdings. Ottawa:
  CIHI.
\item
  Government of Nova Scotia. Nova Scotia Open Data Portal --- including
  the LTC placement waitlist (c39g-gsdd) and the LTC facilities
  directory (x76a-axw2). data.novascotia.ca.
\end{enumerate}

\end{document}
